\documentclass[11pt]{article}
\usepackage[margin=1in]{geometry}
\usepackage{amsmath,amssymb,amsthm,mathtools}
\usepackage{enumitem}
\usepackage{microtype}
\input{glyphtounicode}
\pdfgentounicode=1
\pdfglyphtounicode{mu}{03BC}
\pdfglyphtounicode{bardbl}{2016}
\pdfglyphtounicode{parenleftbig}{0028}
\pdfglyphtounicode{parenrightbig}{0029}
\pdfglyphtounicode{parenleftbigg}{0028}
\pdfglyphtounicode{parenrightbigg}{0029}
\pdfglyphtounicode{parenlefttp}{239B}
\pdfglyphtounicode{parenleftbt}{239D}
\pdfglyphtounicode{parenrighttp}{239E}
\pdfglyphtounicode{parenrightbt}{23A0}
\pdfglyphtounicode{circleplusdisplay}{2A01}
\pdfglyphtounicode{summationdisplay}{2211}
\pdfglyphtounicode{productdisplay}{220F}
\usepackage[colorlinks=true,linkcolor=blue,citecolor=blue,urlcolor=blue]{hyperref}
\usepackage{accsupp}
\NewCommandCopy{\OriginalHookrightarrow}{\hookrightarrow}
\NewCommandCopy{\OriginalLongmapsto}{\longmapsto}
\NewCommandCopy{\OriginalUpbracefill}{\upbracefill}
\renewcommand{\hookrightarrow}{\mathrel{
  \BeginAccSupp{method=hex,unicode,ActualText=21AA}
  \OriginalHookrightarrow\EndAccSupp{}}}
\renewcommand{\longmapsto}{\mathrel{
  \BeginAccSupp{method=hex,unicode,ActualText=27FC}
  \OriginalLongmapsto\EndAccSupp{}}}
\renewcommand{\upbracefill}{
  \BeginAccSupp{method=hex,unicode,ActualText=23DF}
  \OriginalUpbracefill\EndAccSupp{}}
\hypersetup{pdftitle={Quadratic stabilization of the canonical Foulkes--Howe map},pdfauthor={Rewritten for readability from the OpenAI manuscript of the same title}}
\newtheorem{theorem}{Theorem}
\newtheorem{corollary}[theorem]{Corollary}
\newtheorem{proposition}[theorem]{Proposition}
\newtheorem{lemma}[theorem]{Lemma}
\newtheorem*{criterion}{Diagonal criterion}
\theoremstyle{remark}
\newtheorem{remark}[theorem]{Remark}
\DeclareMathOperator{\Sym}{Sym}
\DeclareMathOperator{\Av}{Av}
\newcommand{\C}{\mathbb C}
\newcommand{\Sa}{\mathfrak S_a}
\newcommand{\D}[2]{D_{#1,#2}}
\newcommand{\unverified}[1]{\par\smallskip\noindent\fbox{\parbox{0.95\linewidth}{\textbf{Unverified step.} #1}}\par\smallskip}
\title{Quadratic stabilization of the canonical Foulkes--Howe map}
\author{Rewritten for readability from the OpenAI manuscript of the same title}
\date{}
\begin{document}
\maketitle
\begin{abstract}
For every finite-dimensional complex vector space $V$, the canonical
Foulkes--Howe map $\Sym^b(\Sym^a V)\to\Sym^a(\Sym^b V)$ is surjective
whenever $a\ge2$ and $b\ge a(a-1)$; the bound does not depend on $\dim V$.
This is a step-by-step rewrite of the source's proof, with statements,
constants and strategy unchanged.
\end{abstract}

\section{The result}\label{sec:introduction}

\paragraph{Setting.}
Let $V$ be a finite-dimensional complex vector space and
$\Sym V=\bigoplus_{j\ge0}\Sym^jV$ its symmetric algebra. Products of vectors
are always taken in $\Sym V$. Fix an integer $a\ge1$. The algebra
$\mathcal P=(\Sym V)^{\otimes a}$ is commutative under slotwise
multiplication. (The source calls it $P_a$; it is renamed because $P$ will
denote a product of vectors.) The symmetric group $\Sa$ permutes the tensor slots,
acting by algebra automorphisms. Put
\[
 R_j=\bigl((\Sym^j V)^{\otimes a}\bigr)^{\Sa},\qquad R=\bigoplus_{j\ge0}R_j .
\]
Invariants of algebra automorphisms are closed under multiplication, and
$R_iR_j\subseteq R_{i+j}$, so $R$ is a graded commutative algebra.

\paragraph{Averaging.}
For any complex vector space $W$ there is a linear isomorphism
$\Sym^aW\to(W^{\otimes a})^{\Sa}$,
\[
 w_1\cdots w_a\longmapsto
 \Av(w_1\otimes\cdots\otimes w_a)
 =\frac1{a!}\sum_{\pi\in\Sa}w_{\pi(1)}\otimes\cdots\otimes w_{\pi(a)} .
\]
It is well defined because the right side is symmetric in the $w_i$. Its
inverse is the quotient projection $W^{\otimes a}\to\Sym^aW$ restricted to
invariant tensors (standard, since $a!\ne0$). With $W=V$ and $W=\Sym^bV$
this identifies $R_1=\Sym^aV$ and $R_b=\Sym^a(\Sym^bV)$.

\paragraph{The map.}
For $b\ge0$, multiplication in $R$ gives a linear map $\Sym^bR_1\to R_b$,
$z_1\cdots z_b\mapsto z_1\cdots z_b$. Under the identifications above this is
the \emph{canonical Foulkes--Howe map}
\begin{equation}\label{eq:map}
 \mu_{a,b,V}:\Sym^b(\Sym^a V)\longrightarrow\Sym^a(\Sym^b V).
\end{equation}
Its image is $A_b$, the span of all products of $b$ elements of $R_1$. This
is the degree-$b$ part of the subalgebra $A\subseteq R$ generated by $R_1$.

\begin{theorem}\label{thm:main}
For every integer $a\ge2$, every integer $b\ge a(a-1)$, and every
finite-dimensional complex vector space $V$, the map
$\mu_{a,b,V}$ in \eqref{eq:map} is surjective. Equivalently, $A_b=R_b$.
\end{theorem}

For $a=1$ the map is the identity in every degree.

\paragraph{Proof at a glance.}
\emph{Key idea: dualize.} Section~\ref{sec:map} shows that $\mu_{a,b,V}$ is
surjective if and only if every symmetric $a$-linear form $T$ on $\Sym^bV$
with $T(P,\ldots,P)=0$ for all products $P=u_1\cdots u_b$ of $b$ vectors is
zero. Proposition~\ref{prop:vanishing} proves that such a $T$ is zero. The
proof is by induction on $a$, with $r=a-1$ and $m=b-r$, in four steps.
\begin{enumerate}[leftmargin=2em,itemsep=1pt,topsep=3pt]
\item \emph{Binary form.} Since $x^r+ct^r$ factors into linear factors,
$(x^r+ct^r)Q$ is a product of $b$ vectors for every product $Q$ of $m$
vectors. The coefficient of $c$ gives $T(x^rQ,\ldots,x^rQ,t^rQ)=0$.
\item \emph{Freeing the last slot.} Starting from $T(x^rQ,\ldots,x^rQ,t^b)$,
$m$ derivative operators each trade one factor $t$ for a factor $y_i$ of
$Q$ and arrive at the expression of Step~1. Lemma~\ref{lem:shift} says each operator is injective, so
$T(x^rQ,\ldots,x^rQ,t^b)=0$.
\item \emph{Induction.} For fixed $x,t$, the form
$(Q_1,\ldots,Q_r)\mapsto T(x^rQ_1,\ldots,x^rQ_r,t^b)$ is symmetric and
$r$-linear on $\Sym^mV$, and by Step~2 its diagonal vanishes on products. The induction hypothesis shows
it is zero.
\item \emph{Removing $x^r$.} Starting from $T(v_1^b,\ldots,v_r^b,t^b)$,
$r^2$ operators each trade one factor $v_i$ for a factor of one shared
variable $x$ and arrive at an expression that is zero by Step~3. Lemma~\ref{lem:shift} again gives
injectivity, so $T$ vanishes on pure powers, which span.
\end{enumerate}
\emph{Dependencies.} Lemma~\ref{lem:shift} feeds Steps~2 and~4. Each step
uses the previous one. The spanning fact (2.1) feeds the criterion, the base
case and the end of Step~4. \emph{Where the bound enters.} The hypothesis
$b\ge a(a-1)=r(r+1)$ is $m\ge r^2$. Step~3 uses it as $m\ge r(r-1)$. Step~4
uses it as $m+1>r^2-1$: the $r^2$ operators there all raise the degree of
the same variable $x$, which is why the bound is quadratic.

\paragraph{Context.}
This paragraph reports the source's account of the literature. It is not
used in the proof and has not been re-derived here. In positive dimension,
Raicu--Sam--Weyman \cite[Section~2]{RSW} formulate the map as
$\Sym^b(D^aV)\to D^a(\Sym^bV)$ with divided powers realized as invariant
tensors, and averaging transports it to \eqref{eq:map}. The ring $A$ is the
coordinate ring of the affine Chow variety of products of $a$ linear forms
in $\Sym^a(V^*)$, with normalization $R$
\cite[Section~2, Theorem~2.3]{RSW}; see also
\cite[Section~7.6, Propositions~7.11--7.12]{Landsberg}. Brion proved
surjectivity for all large $b$ \cite{Brion1993} and gave an effective bound
involving both $a$ and $\dim V$ \cite[Theorem~3.3]{Brion}.
\cite[Theorem~3.9]{RSW} bounds the regularity of related modules, including
the normalization. It does not compute the top nonzero degree of $R/A$, and
that degree determines the least $b_0$ with $\mu_{a,b,V}$ surjective for all
$b\ge b_0$. The source states that Theorem~\ref{thm:main} answers positively the request
for a polynomial bound in \cite[Problem~7.19]{Landsberg}: in that notation
$d=b$, $n=a$, $w=\dim V$, and the bound is $d\ge n(n-1)$, including when
$w=n$. Sharpness is not claimed. The Foulkes conjecture
\cite{Foulkes1950} asks, for $a\le b$, only for \emph{some}
$\mathrm{GL}(V)$-equivariant injection
$\Sym^a(\Sym^b V)\hookrightarrow\Sym^b(\Sym^a V)$
\cite[Conjecture~2.10]{RSW}. Corollary~\ref{cor:comparison} supplies one in
our range. The existence of an embedding does not specify the canonical
map: at $(a,b)=(5,5)$ and $(6,6)$, outside our range, the canonical map has
nonzero kernel in suitable dimensions
\cite[arXiv version, Theorem~6(b),(c) and Remark~7]{CIM}. Successive weight
shifts like those of Steps~2 and~4 occur in Ikenmeyer's proof
\cite[Section~5]{Ikenmeyer} of McKay's propagation theorem \cite{McKay}.
That theorem concerns the canonical maps in the opposite direction with
$\dim V\ge b$: for fixed $a$ and $B\ge a$, injectivity at $(a,B)$
propagates to every $(a,b)$ with $b\ge B$
\cite[Theorem~1.3 and Section~3]{Ikenmeyer}.

\section{The diagonal criterion}\label{sec:map}

Fix integers $a,b\ge1$ and put $E=\Sym^bV$. A \emph{product of $b$ vectors}
is an element $u_1\cdots u_b\in E$ with $u_i\in V$. Repeats and zeros are
allowed.

\paragraph{(2.1) Pure powers span.}
\emph{For $k\ge1$, the powers $u^k$ with $u\in V$ span $\Sym^kV$.} This is
used in (2.2), in the base case of Proposition~\ref{prop:vanishing} and at
the end of its Step~4. It follows from the polarization identity
\[
 k!\,u_1\cdots u_k
 =\sum_{J\subseteq\{1,\ldots,k\}}(-1)^{k-|J|}
   \Bigl(\sum_{j\in J}u_j\Bigr)^k .
\]
To check the identity, expand the right side with the $u_i$ as formal
commuting variables. A monomial whose set of occurring indices is $S$
appears in the $J$-term exactly when $J\supseteq S$, each time with the same
multinomial coefficient. Its total coefficient is that multinomial times
$\sum_{J\supseteq S}(-1)^{k-|J|}=(1-1)^{k-|S|}$, which is zero unless
$S=\{1,\ldots,k\}$. The only degree-$k$ monomial with full index set is
$u_1\cdots u_k$, with multinomial coefficient $k!$. Now substitute vectors.
Products of $k$ vectors span $\Sym^kV$ and $k!\ne0$, so powers span.

\paragraph{(2.2) The image.}
Averaging sends $u^a$ to $z(u):=u^{\otimes a}$, since all $a!$ summands are
equal. By (2.1) with $k=a$, the elements $z(u)$ span $R_1$. Slotwise
multiplication gives
$z(u_1)\cdots z(u_b)=(u_1\cdots u_b)^{\otimes a}$. Products of $b$
elements from a spanning set of $R_1$ span $A_b$, so
\[
 A_b=\operatorname{span}_{\C}
 \{P^{\otimes a}: P\text{ a product of }b\text{ vectors}\}.
\]

\paragraph{(2.3) Functionals and forms.}
For $\lambda\in R_b^*$ define
$T_\lambda(w_1,\ldots,w_a)
 =\lambda\bigl(\Av(w_1\otimes\cdots\otimes w_a)\bigr)$ for $w_i\in E$.
This is a symmetric $a$-linear form on $E$. \emph{Claim:
$\lambda\mapsto T_\lambda$ is a bijection from $R_b^*$ onto the symmetric
$a$-linear forms on $E$.}

\emph{Onto.} Given a symmetric $T$, let $\tilde T:E^{\otimes a}\to\C$ be the
linear map with $\tilde T(w_1\otimes\cdots\otimes w_a)=T(w_1,\ldots,w_a)$,
and let $\lambda$ be its restriction to the subspace
$R_b=(E^{\otimes a})^{\Sa}$. Then
$T_\lambda(w_1,\ldots,w_a)=\frac1{a!}\sum_\pi
T(w_{\pi(1)},\ldots,w_{\pi(a)})=T(w_1,\ldots,w_a)$ by symmetry.

\emph{One-to-one.} Simple tensors span $E^{\otimes a}$ and $\Av$ fixes
invariant tensors, so the elements $\Av(w_1\otimes\cdots\otimes w_a)$ span
$R_b$. Hence $T_\lambda=0$ forces $\lambda=0$.

\paragraph{(2.4) Conclusion.}
For a product $P$ we have $\Av(P^{\otimes a})=P^{\otimes a}$, so
$T_\lambda(P,\ldots,P)=\lambda(P^{\otimes a})$. By (2.2), $\lambda$
vanishes on $A_b$ exactly when the diagonal of $T_\lambda$ vanishes on all
products of $b$ vectors. Since $R_b$ is finite-dimensional, $A_b=R_b$ if and only if
$0$ is the only functional vanishing on $A_b$. With (2.3) this gives the
following.

\begin{criterion}
The map $\mu_{a,b,V}$ is surjective if and only if every symmetric
$a$-linear form $T$ on $\Sym^b V$ satisfying $T(P,\ldots,P)=0$ for all
products $P$ of $b$ vectors is zero.
\end{criterion}

\section{Differential injectivity}\label{sec:differential}

Assume $n=\dim V>0$ and fix a basis $e_1,\ldots,e_n$ of $V$. A \emph{vector
variable} is a block $U=(U_1,\ldots,U_n)$ of independent scalar variables.
We write $u=\sum_kU_ke_k$ for the vector it parametrizes and, in
Section~\ref{sec:vanishing}, use the lower-case letter for both. For two
different blocks $U,Y$ define
\[
 \D{Y}{U}=\sum_{k=1}^n Y_k\frac{\partial}{\partial U_k}.
\]
It lowers the $U$-degree by one and raises the $Y$-degree by one.
Section~\ref{sec:vanishing} uses two elementary facts.
\begin{itemize}[leftmargin=2.6em,itemsep=1pt,topsep=3pt]
\item[(D1)] A polynomial in finitely many variables over $\C$ that vanishes
at every point is the zero polynomial. This follows by induction on the
number of variables from the univariate root bound.
\item[(D2)] Let $T$ be a multilinear form on $\Sym^bV$, let $j\ge1$, and
let one slot be $u^jw$, where $w$ and all other slots do not involve $U$.
Then
$\D{Y}{U}\,T(\ldots,u^jw,\ldots)=j\,T(\ldots,u^{j-1}y\,w,\ldots)$.
Indeed, the Leibniz rule in the commutative algebra $\Sym V$ gives
$\partial u^j/\partial U_k=j\,u^{j-1}e_k$. Multiply by $Y_k$, sum over
$k$, and use linearity of $T$ in that slot.
\end{itemize}

The next lemma is what allows a derivative to be cancelled. It is used in
Steps~2 and~4 of Proposition~\ref{prop:vanishing}. The source compares it
with positive-weight injectivity in $\mathfrak{sl}_2$ weight strings
\cite[Section~11.1, Claim~11.4 and equation~(11.5)]{FultonHarris} and with
\cite[Section~5, Claim~5.1]{Ikenmeyer}. The proof below uses neither.

\begin{lemma}\label{lem:shift}
Let $p,q$ be nonnegative integers. On polynomials homogeneous of degree $p$
in $U$ and degree $q$ in $Y$, the operator $\D{Y}{U}$ is injective if
$p>q$. The polynomials may depend on any finite collection of additional
variables independent of $U$ and $Y$.
\end{lemma}

\begin{proof}
Write $D=\D{Y}{U}$ and $D'=\D{U}{Y}$. (The source calls the latter $E$; it is renamed
to avoid a clash with $E=\Sym^bV$.) Let $h$ have bidegree $(p,q)$.

\emph{Step 1: inner product.} Let $X$ list all scalar variables, including
the additional ones. Give the polynomial ring the Fischer (apolar)
Hermitian inner product
\cite[Section~2, equation~(8) and Proposition~3]{AldazRender}: the
monomials are orthonormal, $\|X^\alpha\|^2=1$.
It is positive definite because it has an orthonormal basis.

\emph{Step 2: adjoints.} For each coordinate $X_k$,
$\langle\partial_{X_k}f,g\rangle=\langle f,X_kg\rangle$. It suffices to
take $f=X^\alpha$ and $g=X^\beta$. Both sides are then $1$ if
$\alpha=\beta+e_k$ and $0$ otherwise, because
the monomials are orthonormal. Hence
$D^*=\sum_k(\partial_{U_k})^*(Y_k)^*=\sum_kU_k\partial_{Y_k}=D'$, and
likewise $(D')^*=D$.

\emph{Step 3: commutator.} Using $[\partial_{Y_j},Y_k]=\delta_{jk}$ and
$[U_j,\partial_{U_k}]=-\delta_{jk}$, with all other pairs commuting,
\[
 [U_j\partial_{Y_j},\,Y_k\partial_{U_k}]
 =U_j[\partial_{Y_j},Y_k]\partial_{U_k}+Y_k[U_j,\partial_{U_k}]\partial_{Y_j}
 =\delta_{jk}\bigl(U_j\partial_{U_j}-Y_j\partial_{Y_j}\bigr).
\]
Summing over $j,k$ gives $D'D-DD'=N_U-N_Y$, where
$N_U=\sum_kU_k\partial_{U_k}$ and $N_Y=\sum_kY_k\partial_{Y_k}$. These
Euler operators multiply each monomial by its $U$-degree and $Y$-degree
respectively, so $N_Uh=ph$ and $N_Yh=qh$.

\emph{Step 4: norm identity.} By Steps~2 and~3,
\begin{equation}\label{eq:norm}
 \|Dh\|^2-\|D'h\|^2
 =\langle D'Dh,h\rangle-\langle DD'h,h\rangle
 =\langle(N_U-N_Y)h,h\rangle=(p-q)\|h\|^2.
\end{equation}

\emph{Step 5.} If $Dh=0$ and $p>q$, the left side of \eqref{eq:norm} is
$\le0$ and the right side is $\ge0$. So $(p-q)\|h\|^2=0$ and $h=0$.
\end{proof}

\emph{How it is applied.} To deduce $g=0$ from $D_k\cdots D_1g=0$, peel
operators off from the outside. This requires the degree inequality of the
lemma for the bidegree of $D_{i-1}\cdots D_1g$, for each $i$.

\section{Vanishing on products}\label{sec:vanishing}

\begin{proposition}\label{prop:vanishing}
Let $V$ be a finite-dimensional complex vector space, and let $a\ge1$ and
$b\ge1$ be integers satisfying $b\ge a(a-1)$.
If a symmetric $a$-linear form $T$ on $\Sym^b V$ satisfies
\begin{equation}\label{eq:diagonal}
 T(P,\ldots,P)=0
 \qquad\text{for every product $P$ of $b$ vectors},
\end{equation}
then $T=0$.
\end{proposition}

\begin{proof}
\emph{Plan.} We induct on $a$, proving the statement for all admissible $b$
and all $V$ at once. Steps~1--4 show that $T$ vanishes on all tuples of pure
powers. Then (2.1) finishes the proof.

If $V=0$ then $\Sym^bV=0$ because $b\ge1$, so $T=0$. Assume $\dim V>0$ and
use the conventions of Section~\ref{sec:differential}.

\emph{Base case $a=1$.} $T$ is a linear form vanishing on products of $b$
vectors, which span $\Sym^bV$.

\emph{Inductive step.} Let $a\ge2$ and assume the statement for $a-1$. Put
\[
 r=a-1\ge1,\qquad m=b-r .
\]
The hypothesis $b\ge a(a-1)=r(r+1)$ gives $m\ge r^2\ge1$.

\paragraph{Step 1: a split binary form.}
\emph{Claim: for all $x,t\in V$ and every product $Q$ of $m$ vectors,}
\begin{equation}\label{eq:binary}
 T(\underbrace{x^rQ,\ldots,x^rQ}_{r\text{ slots}},t^rQ)=0.
\end{equation}

(a) \emph{For every $c\in\C$, $(x^r+ct^r)Q$ is a product of $b$ vectors.}
Choose $\alpha\in\C$ with $\alpha^r=-c$ and let $\zeta=e^{2\pi i/r}$.
Homogenizing $z^r-1=\prod_j(z-\zeta^j)$ gives, in $\C[X,Y]$,
\[
 X^r+cY^r=\prod_{j=0}^{r-1}(X-\alpha\zeta^jY).
\]
For $c=0$ this reads $X^r=X^r$. Since $\Sym V$ is commutative, there is an
algebra homomorphism $\C[X,Y]\to\Sym V$ with $X\mapsto x$, $Y\mapsto t$.
This needs no independence of $x$ and $t$. Applying it,
$x^r+ct^r=\prod_j(x-\alpha\zeta^jt)$ is a product of $r$ vectors, and
multiplying by $Q$ gives a product of $r+m=b$ vectors.

(b) \emph{Coefficient of $c$.} By (a) and \eqref{eq:diagonal},
$T(P_c,\ldots,P_c)=0$ for all $c$, where $P_c=x^rQ+c\,t^rQ$. By
multilinearity and symmetry,
\[
 T(P_c,\ldots,P_c)=\sum_{k=0}^{a}\binom ak c^k\,
 T(\underbrace{x^rQ,\ldots}_{a-k},\underbrace{t^rQ,\ldots}_{k}).
\]
A polynomial in $c$ vanishing for all $c$ has zero coefficients. The
coefficient of $c$ is $a\,T(x^rQ,\ldots,x^rQ,t^rQ)$ with $a-1=r$ copies of
$x^rQ$, and $a\ne0$. This proves \eqref{eq:binary}.

\paragraph{Step 2: freeing the last slot.}
\emph{Claim: for all $x,t\in V$ and every product $Q$ of $m$ vectors,}
\begin{equation}\label{eq:free}
 T(\underbrace{x^rQ,\ldots,x^rQ}_{r\text{ slots}},t^b)=0.
\end{equation}
Fix $x\in V$. Let $t,y_1,\ldots,y_m$ be independent vector variables, put
$Q=y_1\cdots y_m$, and for $0\le k\le m$ define the polynomial
\[
 G_k=T(x^rQ,\ldots,x^rQ,\;t^{b-k}y_1\cdots y_k).
\]
The source's $H_x$ is $G_0$, the left side of \eqref{eq:free}.

(a) \emph{$G_m=0$ as a polynomial.} Since $b-m=r$, the last slot of $G_m$
is $t^rQ$. So $G_m$ vanishes at every point by \eqref{eq:binary}, and (D1)
applies.

(b) \emph{$\D{y_k}{t}G_{k-1}=(b-k+1)\,G_k$ for $1\le k\le m$.} Only the
last slot involves $t$. Apply (D2) with $j=b-k+1$ and
$w=y_1\cdots y_{k-1}$. The scalar $b-k+1$ is positive. Composing gives the
source's formula
$\D{y_k}{t}\cdots\D{y_1}{t}G_0=\frac{b!}{(b-k)!}G_k$.

(c) \emph{Bidegree.} $G_{k-1}$ has degree $p=b-k+1$ in $t$, which occurs
only in the last slot. It has degree $q=r$ in $y_k$: each of the first $r$
slots contains $Q$, which is linear in $y_k$, and the last slot does not
yet contain $y_k$. Thus $p-q=m-k+1\ge1$.

(d) \emph{Descent.} If $G_k=0$, then $\D{y_k}{t}G_{k-1}=0$ by (b), so
$G_{k-1}=0$ by Lemma~\ref{lem:shift} and (c), with the other variables as
additional variables. Starting from (a) and taking $k=m,m-1,\ldots,1$
yields $G_0=0$. Evaluating at arbitrary vectors gives \eqref{eq:free}.

\paragraph{Step 3: induction on the remaining slots.}
\emph{Claim: for all $x,t\in V$,}
\begin{equation}\label{eq:inductive-vanishing}
 T(x^rQ_1,\ldots,x^rQ_r,t^b)=0
 \qquad(Q_1,\ldots,Q_r\in\Sym^m V).
\end{equation}
Fix $x,t$ and let $S_{x,t}(Q_1,\ldots,Q_r)$ denote the left side. It is
$r$-linear because $Q\mapsto x^rQ$ is a linear map $\Sym^mV\to\Sym^bV$, and
symmetric because $T$ is. By \eqref{eq:free}, its diagonal vanishes at
every product of $m$ vectors. The induction hypothesis, for $r$ slots and degree
$m$, requires $m\ge1$ and $m\ge r(r-1)$. Both follow from $m\ge r^2\ge1$.
Hence $S_{x,t}=0$.

\paragraph{Step 4: removing the common power.}
\emph{Claim: $T(v_1^b,\ldots,v_r^b,t^b)=0$ for all
$v_1,\ldots,v_r,t\in V$.}
Fix $t\in V$. Let $v_1,\ldots,v_r,x$ be independent vector variables. For
$s=(s_1,\ldots,s_r)$ with $0\le s_i\le r$ put $|s|=\sum_is_i$ and
\[
 F_s=T(x^{s_1}v_1^{b-s_1},\ldots,x^{s_r}v_r^{b-s_r},\;t^b).
\]
The source's $f$ is $F_{(0,\ldots,0)}$.

(a) \emph{$F_{(r,\ldots,r)}=0$ as a polynomial.} It equals
$T(x^rv_1^m,\ldots,x^rv_r^m,t^b)$. This vanishes at every point by
\eqref{eq:inductive-vanishing} with $Q_i=v_i^m$, and (D1) applies.

(b) \emph{If $s_i<r$ then $\D{x}{v_i}F_s=(b-s_i)\,F_{s+e_i}$}, where
$e_i$ is the $i$-th unit vector. Only slot $i$ involves $v_i$. Apply (D2)
with $j=b-s_i$ and $w=x^{s_i}$. The scalar $b-s_i\ge m+1$ is positive.
Applying $\D{x}{v_i}^{\,r}$ for $i=1,\ldots,r$ in turn gives the source's
formula
$\D{x}{v_r}^{\,r}\cdots\D{x}{v_1}^{\,r}F_{(0,\ldots,0)}
=\bigl(\tfrac{b!}{m!}\bigr)^rF_{(r,\ldots,r)}$.

(c) \emph{Bidegree.} $F_s$ has degree $p=b-s_i$ in $v_i$ and degree
$q=|s|$ in $x$. All operators raise the degree of the same variable $x$.
If $s_i\le r-1$, then $|s|\le r^2-1$ and
\[
 p=b-s_i\ \ge\ b-r+1=m+1\ \ge\ r^2+1\ >\ r^2-1\ \ge\ |s|=q .
\]

(d) \emph{Descent.} If $s_i<r$ and $F_{s+e_i}=0$, then
$\D{x}{v_i}F_s=0$ by (b), so $F_s=0$ by Lemma~\ref{lem:shift} and (c).
Start from $(r,\ldots,r)$ by (a) and lower one coordinate at a time, in
$r^2$ steps that undo the source's composition from the outside. This
gives $F_{(0,\ldots,0)}=0$. Evaluating at arbitrary vectors, for each $t$,
proves the claim.

\paragraph{Conclusion.}
By (2.1) with $k=b$, pure powers span $\Sym^bV$. By Step~4 and
multilinearity, $T=0$. This completes the induction.
\end{proof}

\begin{proof}[Proof of Theorem~\ref{thm:main}]
For $a\ge2$ and $b\ge a(a-1)$ we have $b\ge1$.
Proposition~\ref{prop:vanishing} then gives the vanishing statement in the
diagonal criterion of Section~\ref{sec:map}, and the criterion gives
surjectivity.
\end{proof}

\section{Representation comparison and scope}\label{sec:comparison}

\begin{corollary}\label{cor:comparison}
For every finite-dimensional complex vector space $V$, every integer
$a\ge2$, and every integer $b\ge a(a-1)$, there exists a
$\mathrm{GL}(V)$-equivariant embedding
\[
 \Sym^a(\Sym^b V)\hookrightarrow\Sym^b(\Sym^a V).
\]
In particular, this gives the sixth symmetric-power comparison for $b\ge30$.
\end{corollary}

The source does not define the phrase in the last sentence. We read it as
the case $a=6$, where $a(a-1)=30$: an embedding
$\Sym^6(\Sym^bV)\hookrightarrow\Sym^b(\Sym^6V)$ for $b\ge30$.

\begin{proof}
$\mathrm{GL}(V)$ acts on $\Sym V$ by algebra automorphisms, hence on
$\mathcal P$, commuting with $\Sa$. So averaging and multiplication are
equivariant, and the surjection $\mu_{a,b,V}$ of Theorem~\ref{thm:main} is
equivariant. Finite-dimensional polynomial representations over $\C$ are
completely reducible
\cite[Chapter~I, Appendix~A, Sections~7--8]{Macdonald}. Choose an
invariant complement $M$ to $\ker\mu_{a,b,V}$. The restriction
$\mu_{a,b,V}|_M$ is an equivariant isomorphism onto $\Sym^a(\Sym^bV)$. Its
inverse followed by the inclusion $M\subseteq\Sym^b(\Sym^aV)$ is the
required embedding. This shows existence only. It does not identify the
embedding with a prescribed canonical map.
\end{proof}

\begin{remark}
Nothing in the proof depends on $\dim V$. Lemma~\ref{lem:shift} holds for
any number of coordinates and additional variables, and every degree
comparison involves only $a$ and $b$. So no reduction to a particular
dimension is needed. The proof uses positivity of the Fischer inner
product, splitting of binary forms over $\C$, and the nonvanishing in
characteristic zero of the scalars $k!$, $a$, $b!/(b-k)!$ and
$(b!/m!)^r$.
\end{remark}

\section*{Verification checklist}

\begin{enumerate}[leftmargin=2em,itemsep=1pt,topsep=3pt]
\item \emph{Dualization.} The image is spanned by the $P^{\otimes a}$
(2.2), and $\lambda\mapsto T_\lambda$ is a bijection with
$T_\lambda(P,\ldots,P)=\lambda(P^{\otimes a})$ (2.3, 2.4).
\item \emph{Lemma~\ref{lem:shift}.} Check the adjointness $D^*=D'$
(Step~2), the commutator $D'D-DD'=N_U-N_Y$ (Step~3), and that $h$ is
homogeneous separately in $U$ and in $Y$ wherever the lemma is invoked.
\item \emph{Step~1(a).} $(x^r+ct^r)Q$ really is a product of $b$ vectors,
for dependent $x,t$ and for $c=0$ as well. In Step~1(b) the coefficient of
$c$ has exactly $r$ copies of $x^rQ$.
\item \emph{Step~2(c).} Before the $k$-th operator the $y_k$-degree is
still $r$, not $r+1$, and the $t$-degree is $b-k+1$, so $p-q=m-k+1\ge1$.
\item \emph{Step~3.} The induction hypothesis is applied to
$(a-1,m,V)$. Its requirements $m\ge1$ and $m\ge r(r-1)$ hold.
\item \emph{Step~4(c).} With one shared variable $x$, the chain
$b-s_i\ge m+1\ge r^2+1>r^2-1\ge|s|$ holds before each of the $r^2$
operators. Together with Step~3, this is where $b\ge a(a-1)$ is used.
\item \emph{Passing between values and polynomials.} (D1) is used in
Steps~2(a) and~4(a), and the descents then return to values at arbitrary
vectors.
\end{enumerate}

\begin{thebibliography}{99}
\bibitem{AldazRender}
J.~M.~Aldaz and H.~Render,
\href{https://doi.org/10.1007/s13324-023-00844-4}{\emph{A Fischer type decomposition theorem from the apolar inner product}},
Analysis and Mathematical Physics \textbf{13} (2023), article~91.

\bibitem{Brion1993}
M.~Brion,
\href{https://doi.org/10.1007/BF03026558}{\emph{Stable properties of plethysm: on two conjectures of Foulkes}},
Manuscripta Mathematica \textbf{80} (1993), no.~4, 347--371.

\bibitem{Brion}
M.~Brion,
\href{https://smf.emath.fr/publications/sur-certains-modules-gradues-associes-aux-produits-symetriques}{\emph{Sur certains modules gradu\'es associ\'es aux produits sym\'etriques}},
in \emph{Alg\`ebre non commutative, groupes quantiques et invariants}
(Reims, 1995), S\'eminaires et Congr\`es \textbf{2},
Soci\'et\'e Math\'ematique de France, Paris, 1997, 157--183.

\bibitem{CIM}
M.-W.~Cheung, C.~Ikenmeyer, and S.~Mkrtchyan,
\href{https://doi.org/10.1016/j.jsc.2016.09.002}{\emph{Symmetrizing tableaux and the 5th case of the Foulkes conjecture}},
Journal of Symbolic Computation \textbf{80} (2017), part~3, 833--843.
\href{https://arxiv.org/abs/1509.03944v1}{arXiv:1509.03944v1}.

\bibitem{Foulkes1950}
H.~O.~Foulkes,
\href{https://doi.org/10.1112/jlms/s1-25.3.205}{\emph{Concomitants of the quintic and sextic up to degree four in the coefficients of the ground form}},
Journal of the London Mathematical Society \textbf{25} (1950), no.~3,
205--209.

\bibitem{FultonHarris}
W.~Fulton and J.~Harris,
\href{https://doi.org/10.1007/978-1-4612-0979-9}{\emph{Representation Theory: A First Course}},
Graduate Texts in Mathematics \textbf{129}, Springer-Verlag, New York, 1991.

\bibitem{Ikenmeyer}
C.~Ikenmeyer,
\emph{On McKay's propagation theorem for the Foulkes conjecture},
\href{https://arxiv.org/abs/1509.04957v1}{arXiv:1509.04957v1}, 2015.

\bibitem{Landsberg}
J.~M.~Landsberg,
\href{https://doi.org/10.1007/s11565-014-0202-7}{\emph{Geometric complexity theory: an introduction for geometers}},
Annali dell'Universit\`a di Ferrara \textbf{61} (2015), no.~1, 65--117.
\href{https://arxiv.org/abs/1305.7387}{arXiv:1305.7387}.
Locators refer to the published version.

\bibitem{Macdonald}
I.~G.~Macdonald,
\href{https://doi.org/10.1093/oso/9780198534891.001.0001}{\emph{Symmetric Functions and Hall Polynomials}},
second edition, Oxford Mathematical Monographs, Clarendon Press, Oxford, 1995.

\bibitem{McKay}
T.~McKay,
\href{https://doi.org/10.1016/j.jalgebra.2007.12.003}{\emph{On plethysm conjectures of Stanley and Foulkes}},
Journal of Algebra \textbf{319} (2008), no.~5, 2050--2071.

\bibitem{RSW}
C.~Raicu, S.~V.~Sam, and J.~Weyman,
\href{https://doi.org/10.1007/s10013-021-00527-2}{\emph{On some modules supported in the Chow variety}},
Vietnam Journal of Mathematics \textbf{50} (2022), no.~2, 501--521.
\href{https://arxiv.org/abs/2108.10910v1}{arXiv:2108.10910v1}.
Locators refer to this preprint version.
\end{thebibliography}

\end{document}
